\documentclass[10pt,twocolumn]{article}
\usepackage[margin=0.8in]{geometry}
\usepackage{booktabs}
\usepackage{amsmath}
\usepackage{xurl}
\usepackage[hidelinks]{hyperref}
\usepackage{xcolor}
\usepackage{enumitem}
\setlist{nosep,leftmargin=*}

\title{The Schema Tax Is a Schema Design Tax:\\When Constrained JSON Decoding Costs Accuracy, and When It Helps}
\author{Ali Ashter\\\texttt{hello@ashter.dev}}
\date{October 2026}

\begin{document}
\maketitle

\begin{abstract}
Production systems increasingly force language models to answer in JSON, using constrained decoding to guarantee the
output matches a schema. Earlier work reported that format restrictions hurt reasoning, and others disputed it. We
separate the parts of the question that usually travel together: whether the schema is enforced, whether it is shown
to the model, and in what order its fields come. Across three Qwen3 models (1.7B, 4B, 8B) on GSM8K, AG News and
CoNLL-2003, with 12,600 requests in total, enforcing a schema that is also shown in the prompt costs nothing
measurable on GSM8K at 4B and 8B ($-0.3$ and $-2.7$ points against free-text chain of thought, both within the 95\%
interval) and helps on classification and extraction. The losses come from schema design. Putting the
\texttt{answer} field before the \texttt{reasoning} field cuts GSM8K accuracy by 63 to 69 points, to the level of
answering with no reasoning at all, replicating an earlier finding on a hosted model; 55 of the 56 answers that hit
the length limit were repetition loops. Enforcing a schema the model has not
seen costs 9 to 15 points on GSM8K. Unconstrained JSON is the least reliable option for small models: the 1.7B model
produced unparseable output for half of its answers, in part by copying the JSON Schema itself. Showing the schema is
not always better: on extraction it makes 4B and 8B models over-fill a catch-all field, and hiding it raises F1 by
10 to 11 points. Code, prompts and every response are public at \url{https://github.com/alleyxd/SchemaTax}.
\end{abstract}

\section{Introduction}

Structured output is now a default feature of model serving. OpenAI-compatible servers such as vLLM \cite{vllm} accept
a JSON Schema and use a grammar engine \cite{outlines,xgrammar} to mask, at every step, the tokens that would break it.
The output is guaranteed to parse, which removes a whole class of production failures.

Whether this costs accuracy has been argued in both directions. Tam et al.\ \cite{speakfreely} found that format
restrictions degrade reasoning, and a widely read response \cite{saywhat} argued that the effect came from prompt and
parsing choices rather than from constraints. Much of the disagreement comes from changing several things at once.
A typical ``JSON mode'' setup differs from free-text prompting in at least three ways: the decoder is constrained,
the prompt describes the format, and the format fixes an order in which the model must commit to its answer.

We vary these separately on three local models and three tasks:

\begin{itemize}
\item \textbf{natural}: the task's usual free-text format, parsed with a regular expression;
\item \textbf{json\_free}: JSON requested, the schema shown in the prompt, decoding unconstrained;
\item \textbf{json\_constrained}: the schema shown and enforced;
\item \textbf{json\_blind}: the schema enforced but not shown (the prompt only says ``Respond in JSON'');
\item on GSM8K, two field orders under enforcement: \textbf{answer\_first} and \textbf{answer\_only}.
\end{itemize}

Our findings, in short: enforcement itself is nearly free once a model is large enough, and it is a net gain for
small models; the order of fields matters far more than enforcement; and whether to show the schema depends on the
task.

\section{Related work}

Grammar-constrained decoding goes back to work on structured NLP tasks \cite{gcd}, with efficient engines built on
finite-state machines \cite{outlines} and pushdown automata \cite{xgrammar}. JSONSchemaBench \cite{jsonschemabench}
compares engines on coverage, speed and the quality of constrained outputs. Tam et al.\ \cite{speakfreely} reported
that JSON, XML and YAML restrictions lower reasoning accuracy; the response by Kurt \cite{saywhat} attributed most of
that drop to the prompts and the answer parser, and used a reasoning-first schema shown in the prompt. Castillo
\cite{castillo} showed that field order alone is decisive for GPT-4o-mini on GSM8K: 94.2\% with reasoning first,
31.8\% with the answer first. That reasoning must come before the answer is the premise of chain-of-thought
prompting \cite{cot}. We replicate the field-order effect on open models across sizes, and add what these works did
not isolate: schema visibility under enforcement, the failure modes of unconstrained JSON, and how both change with
model size. Label bias in small models, which our classification results touch on, was studied by
Zhao et al.\ \cite{zhao}.

\section{Setup}

\textbf{Models.} Qwen3-1.7B, Qwen3-4B and Qwen3-8B \cite{qwen3} in bf16, served by vLLM 0.29.0 with its default
structured-output backend, thinking disabled, temperature 0, on one RTX 5090.

\textbf{Tasks.} 300 examples each, drawn with a fixed seed: GSM8K \cite{gsm8k} test questions (exact numeric match,
up to 1,024 new tokens), AG News test articles (four-way topic, exact label match) and CoNLL-2003 test sentences with
at least one entity (entity-level F1 over persons, organizations, locations and misc, exact string match).

\textbf{Prompts.} Each task has one base instruction shared by all conditions. The \textbf{natural} condition adds
the usual format: ``finish with a line \emph{The answer is N}'' for GSM8K, ``answer with the label only'' for AG
News, and four labelled lines for CoNLL. The schema conditions add either the full JSON Schema or only ``Respond in
JSON''. The GSM8K schema has a \texttt{reasoning} string followed by an \texttt{answer} number; the AG News schema has a
\texttt{label} enum; the CoNLL schema has one string array per entity type.

\textbf{Statistics.} We report the paired difference to the natural condition on the same examples with a 95\%
bootstrap interval (2,000 resamples).

\section{Results}

\begin{table}[t]
\centering\small
\begin{tabular}{lrrr}
\toprule
GSM8K & 1.7B & 4B & 8B \\
\midrule
natural & 0.827 & 0.923 & 0.927 \\
json\_free & 0.370 & 0.910 & 0.820 \\
json\_constrained & 0.710 & 0.897 & 0.923 \\
json\_blind & 0.697 & 0.770 & 0.840 \\
answer\_first & 0.133 & 0.240 & 0.297 \\
answer\_only & 0.120 & 0.260 & 0.303 \\
\bottomrule
\end{tabular}
\caption{GSM8K accuracy, 300 questions per cell.}
\label{tab:gsm}
\end{table}

\subsection{Enforcement is nearly free when the schema is shown}

On GSM8K (Table~\ref{tab:gsm}), shown-and-enforced JSON is within noise of free-text chain of thought at 8B
($-0.003$, interval $[-0.037, +0.027]$) and 4B ($-0.027$, $[-0.060, +0.007]$). It also uses less than half the
tokens: 128 against 277 at 8B, as the reasoning written inside the JSON string is much terser. The 1.7B model does
lose 11.7 points ($[-0.170, -0.067]$), so a tax exists, but only at the smallest size we tested.

\subsection{Field order matters more than enforcement}

With \texttt{answer} before \texttt{reasoning}, accuracy falls by 63 to 69 points at every size, to the level of the
answer-only schema. The reasoning written after the answer cannot change it, so it is pure cost. Worse, it is
sometimes not reasoning at all: 7, 19 and 30 of these responses (1.7B, 4B, 8B) hit the 1,024-token limit. In 55 of
these 56 cases the model, having committed to a number, repeats the same justification until it is cut off, and in all
56 the committed number was wrong. The size of the drop matches what Castillo \cite{castillo} reported for
GPT-4o-mini.

\subsection{Hidden schemas cost reasoning accuracy}

Enforcing the same reasoning-first schema without showing it costs 8.7, 15.3 and 13.0 points on GSM8K (8B, 4B,
1.7B). The model sees only the keys and types as the grammar forces them, token by token, and is never told that a
\texttt{reasoning} field is coming, so it writes less of it. Blind responses averaged 112 to 126 tokens against 121 to
134 for shown ones.

\subsection{Unconstrained JSON fails on small models}

Asking for JSON without enforcing it is the least reliable choice for small models. At 1.7B only 47\% of GSM8K
answers and 50\% of AG News answers parse. Two failure modes dominate. Models start writing prose with LaTeX math
and then emit JSON containing invalid backslash escapes (31 of the 35 unparseable 8B GSM8K answers, and 118 of 158 at 1.7B). And the 1.7B
model \emph{echoes the schema}: on 297 of 300 CoNLL answers it reproduced the JSON Schema's own structure, with
\texttt{"type": "object"} and its answers nested under \texttt{"properties"}. Unwrapping these echoes raises its AG
News score from 0.410 to 0.820, the same as with enforcement, so the model knows the answer and only gets the format
wrong. Enforcement removes both failure modes.

\begin{table}[t]
\centering\small
\begin{tabular}{llrrr}
\toprule
& & 1.7B & 4B & 8B \\
\midrule
AG News & natural & 0.507 & 0.777 & 0.823 \\
& json\_free & 0.410 & 0.853 & 0.803 \\
& json\_constrained & 0.820 & 0.850 & 0.803 \\
& json\_blind & 0.400 & 0.823 & 0.783 \\
\midrule
CoNLL F1 & natural & 0.321 & 0.476 & 0.380 \\
& json\_free & 0.000 & 0.508 & 0.476 \\
& json\_constrained & 0.430 & 0.511 & 0.474 \\
& json\_blind & 0.387 & 0.620 & 0.574 \\
\bottomrule
\end{tabular}
\caption{Classification accuracy and extraction F1, 300 examples per cell.}
\label{tab:cls}
\end{table}

\subsection{Classification and extraction gain from JSON}

On AG News and CoNLL (Table~\ref{tab:cls}), shown-and-enforced JSON matches or beats the natural format in five of
six cells. The largest gain, $+0.313$ for the 1.7B model on AG News, is mostly a fix for label bias: in the natural
format it answered ``World'' for 212 of 300 articles. Shown the schema with its enum of four labels, it spread its
answers across all four labels (98, 75, 68 and 59); enforced against the same enum without seeing it (blind), it
answered ``World'' 255 times and scored 0.400. The labels are
listed in the instruction in every condition, so it is the explicit enum in the schema, not the list itself, that
makes the difference.

\subsection{Showing the schema can hurt extraction}

CoNLL is the one place where hiding the schema helps at 4B and 8B: F1 rises from 0.511 to 0.620 and from 0.474 to
0.574. Recall is unchanged (0.77 against 0.76 at 4B, 0.73 against 0.75 at 8B); precision rises from 0.38 to 0.58 and
from 0.36 to 0.48. Almost all of the difference is the catch-all \texttt{misc} array. Shown a schema with four arrays,
the models tend to fill each one, putting race times, scores and dates into \texttt{misc}: 627 and 687 false
positives there, against 199 and 393 when the schema is hidden. The 1.7B model behaves the other way (332 shown, 896
blind), so this is not a rule but a reminder that the schema is part of the prompt and can bias it.

\section{Recommendations}

\begin{itemize}
\item Enforce the schema. At 4B and above it costs little or nothing, and below that free-form JSON is too fragile.
\item Put reasoning fields before answer fields. An answer-first schema throws away the reasoning and can loop.
\item Show the schema for reasoning and classification tasks; for extraction with a catch-all field, measure both, or
tighten the field's description.
\item If you parse JSON from a small model without enforcement, expect schema echoes and handle them.
\end{itemize}

\section{Limitations}

We test one model family at three sizes, one serving stack and one grammar backend, with thinking disabled and greedy
decoding. Results with reasoning modes or other families may differ. Each cell has 300 examples, so differences of a
few points are within noise, and we report intervals for that reason. CoNLL F1 uses exact string match, which
penalises boundary differences; absolute scores are low for every condition, and we draw conclusions only from the
differences. We did not tune prompts per condition, deliberately, so a better-worded natural prompt could close some
gaps.

\section{Conclusion}

``Does JSON mode hurt accuracy?'' has a more useful answer than yes or no. Constrained decoding with a shown,
reasoning-first schema is close to free on the models people actually deploy, and it makes small models far more
reliable. The accuracy people lose usually comes from the schema itself: answer fields placed before reasoning, schemas
the model never sees, and catch-all fields it feels obliged to fill.

\begin{thebibliography}{99}\small
\bibitem{vllm} W.~Kwon, Z.~Li, S.~Zhuang, et al. Efficient memory management for large language model serving with PagedAttention. SOSP 2023. arXiv:2309.06180.
\bibitem{outlines} B.~T.~Willard, R.~Louf. Efficient guided generation for large language models. arXiv:2307.09702, 2023.
\bibitem{xgrammar} Y.~Dong, C.~F.~Ruan, Y.~Cai, et al. XGrammar: flexible and efficient structured generation engine for large language models. arXiv:2411.15100, 2024.
\bibitem{gcd} S.~Geng, M.~Josifoski, M.~Peyrard, R.~West. Grammar-constrained decoding for structured NLP tasks without finetuning. EMNLP 2023. arXiv:2305.13971.
\bibitem{jsonschemabench} S.~Geng, H.~Cooper, M.~Moskal, et al. JSONSchemaBench: a rigorous benchmark of structured outputs for language models. arXiv:2501.10868, 2025.
\bibitem{speakfreely} Z.~R.~Tam, C.-K.~Wu, Y.-L.~Tsai, et al. Let me speak freely? A study on the impact of format restrictions on performance of large language models. arXiv:2408.02442, 2024.
\bibitem{saywhat} W.~Kurt. Say what you mean: a response to ``Let Me Speak Freely''. \url{https://blog.dottxt.ai/say-what-you-mean.html}, 2024.
\bibitem{castillo} D.~Castillo. Order of fields in structured output can hurt LLMs output. \url{https://www.dsdev.in/order-of-fields-in-structured-output-can-hurt-llms-output}, 2025.
\bibitem{cot} J.~Wei, X.~Wang, D.~Schuurmans, et al. Chain-of-thought prompting elicits reasoning in large language models. NeurIPS 2022. arXiv:2201.11903.
\bibitem{zhao} T.~Z.~Zhao, E.~Wallace, S.~Feng, D.~Klein, S.~Singh. Calibrate before use: improving few-shot performance of language models. ICML 2021. arXiv:2102.09690.
\bibitem{qwen3} A.~Yang, A.~Li, B.~Yang, et al. Qwen3 technical report. arXiv:2505.09388, 2025.
\bibitem{gsm8k} K.~Cobbe, V.~Kosaraju, M.~Bavarian, et al. Training verifiers to solve math word problems. arXiv:2110.14168, 2021.
\end{thebibliography}

\end{document}
